\documentclass[11pt]{article}
\usepackage[a4paper,margin=1.9cm]{geometry}
\usepackage[T1]{fontenc}
\usepackage{textcomp}
\usepackage[spanish,es-noshorthands]{babel}
\usepackage{amsmath,amssymb}
\usepackage{enumitem}
\usepackage{xcolor}
\usepackage{parskip}
\usepackage{lmodern}
\renewcommand{\familydefault}{\sfdefault}

\definecolor{unalblue}{RGB}{31,73,124}

\newcommand{\examheader}{%
  \noindent
  \begin{minipage}[t]{0.6\linewidth}
    {\small\bfseries MÉTODOS NUMÉRICOS --- Taller 9}\\
    {\small Ejercicio 5 (fragmento)}
  \end{minipage}%
  \hfill
  \begin{minipage}[t]{0.35\linewidth}
    \raggedleft\small Escuela de Matemáticas. Universidad Nacional de Colombia, Sede Medellín.
  \end{minipage}\\[2pt]
  \rule{\linewidth}{0.6pt}
}
\newcommand{\examfooter}{%
  \vfill
  \rule{\linewidth}{0.4pt}\\[1pt]
  {\scriptsize\textit{Transcripción digital --- MathUNAL}\hfill\textcolor{unalblue}{\textbf{\thepage}}}
}
\pagestyle{empty}

\begin{document}
\examheader

\vspace{10pt}

\noindent 5. (25\%) Considere el problema parabólico
\[
\begin{cases}
u_t(x,t) = \pi^2 u_{xx}(x,t), & 0 < x < \dfrac{\pi}{4}, \quad t > 0, \\[4pt]
u(x,0) = f(x), & 0 \le x \le \dfrac{\pi}{4}, \\[4pt]
u(0,t) = g_1(t), & t > 0, \\[4pt]
u\!\left(\dfrac{\pi}{4},t\right) = g_2(t), & t > 0,
\end{cases}
\]
donde $g_1(t)$ y $g_2(t)$ no nulas para valores de $t > 0$.

\vspace{6pt}
\noindent\textbf{Plantée} el sistema de ecuaciones lineales necesario para aproximar la solución $u$ en el tiempo $t = 0.05$, sistema que se obtiene cuando se emplea el método implícito (regresivo o hacia atrás) con tamaños de paso $h = \dfrac{\pi}{8}$ y $k = 0.05$.

\vspace{6pt}
\noindent Tenga en cuenta los siguientes valores conocidos de $u$

\begin{center}
\begin{tabular}{c|ccccc}
 & $x_0$ & $x_1$ & $x_2$ & $x_3$ & $x_4$ \\
\hline
$u(x,0)$ & 0 & 0.1951 & 0.3827 & 0.5556 & 0
\end{tabular}
\end{center}

\vspace{6pt}
\noindent $u(0,0.05) = 0.05$ \quad y \quad $u\!\left(\dfrac{\pi}{4},0.05\right) = 0.025$.

\examfooter
\end{document}
